\documentclass[11pt,a4paper]{article}
\usepackage[margin=24mm]{geometry}
\usepackage[T1]{fontenc}
\usepackage{lmodern,microtype}
\usepackage{amsmath,amssymb,amsthm,mathtools}
\usepackage[hidelinks]{hyperref}
\newtheorem{theorem}{Theorem}
\newtheorem{proposition}[theorem]{Proposition}
\newtheorem{lemma}[theorem]{Lemma}
\newtheorem{corollary}[theorem]{Corollary}
\newtheorem{remark}[theorem]{Remark}
\newcommand{\N}{\mathbb N}
\newcommand{\one}{\mathbf 1}
\newcommand{\calC}{\mathcal C}
\newcommand{\calE}{\mathcal E}
\title{Live Parity False Positives and Exact Return Certificates\\
\large An internal audit of the weighted Collatz clipping-tail strategy}
\author{Research note --- independent review pending}
\date{17 November 2026}
\begin{document}
\maketitle
\begin{abstract}
We audit the proposed improvement of the Collatz clipping-tail bound by
restricting dangerous parity cylinders to trajectories that remain above the
first-passage cutoff during the witnessing prefix.  We give an explicit
infinite family of \emph{real} accelerated $3x+1$ trajectories that satisfies
this survival gate, with an arbitrarily long all-odd prefix, but makes only
two visits to its initial dyadic shell before first passage. Thus the
live-parity gate is still only a necessary, not sufficient, return certificate.
Conditional on a separately stated, presently unverified uniform lower bound
for logarithmic predecessor densities, the same family produces polynomial
spikes in the gated potential.  We then prove an exact return-certificate
identity and the corresponding sharper stopped-potential bound.  Neither
statement proves that the remaining weighted mass is $o(X)$.
\end{abstract}

\section{Set-up and quantifiers}
Let
\[
T(n)=\begin{cases}n/2,&n\text{ even},\\(3n+1)/2,&n\text{ odd},\end{cases}
\qquad \beta=1/200,\quad \gamma=1-2\beta=0.99.
\]
For a large real $X$ let $X_1=X^\gamma$, $W_X=(X^{1-\beta},X]\cap\N$,
$\zeta=\log 2/\log3$, and $\zeta_0=11/20$. Let
\[
 t_X=\left\lceil\frac{3\beta\log X}{(\zeta-\zeta_0)\log3}\right\rceil,
\qquad b_X=\left\lceil(\log\log X)^2\right\rceil.
\]
Write $H_\ell(y)$ for the number of odd steps among the first $\ell$
applications of $T$. The \emph{live parity danger set} is
\[
\calC^{\mathrm{live}}_{X,b}
=\left\{y\in W_X:\ \exists\ell\in[b,t_X],\quad
H_\ell(y)\ge\zeta_0\ell,\quad
T^i(y)>X_1\ \text{for all }0\le i<\ell\right\}.
\]
For a dyadic shell $I\subset W_X$, write $N_I^+(y)$ for the
number of \emph{later} visits of the orbit of $y$ to $I$ before first
hitting $[1,X_1]$. The first-passage stopping time is $\tau_{X_1}(y)$.

\section{A surviving parity witness with only two shell visits}
\begin{theorem}[Live danger does not imply many returns]\label{thm:family}
For each sufficiently large integer $t$, put
\[
 s=\left\lceil(\log t+3)^2\right\rceil,\qquad
 r=\left\lceil t/45\right\rceil.
\]
There exist positive integers $v,n,a,m$ such that
\begin{align*}
 &v<2^{r+2},\quad 3\nmid n,\quad
 n=2^{t+s}v-1,\quad a=2^s v3^t-1,\\
 &m=(v3^{t+s}-1)/2^r\in\N\text{ is odd},\\
 &T^t(n)=a,\quad T^s(a)=2^rm,\quad T^{s+r}(a)=m.
\end{align*}
Set $X=a$. Then, for all sufficiently large $t$:
\begin{enumerate}
\item $m<X^{0.99}$, so the orbit of $a$ reaches the actual
first-passage cutoff;
\item $a\in\calC^{\mathrm{live}}_{X,b_X}$, witnessed by the prefix
$\ell=s$ with $H_s(a)=s$, and $\nu_2(3a+1)=1$;
\item the orbit of $a$ has \emph{exactly two} visits to its own shell
$I=(X/2,X]$ before hitting $[1,X_1]$: the initial point and one
point on the later halving descent.
\end{enumerate}
Thus survival-gated parity danger can hold along an arbitrarily long
prefix even when $N_I^+(a)=1$.
\end{theorem}
\begin{proof}
Because $3^{t+s}$ is odd, it is invertible modulo $2^{r+1}$.
Choose the unique odd residue $v_0\in[1,2^{r+1})$ satisfying
\[
 3^{t+s}v_0\equiv1+2^r\pmod{2^{r+1}}.
\]
Of $v_0$ and $v_0+2^{r+1}$ choose $v$ for which
$3\nmid 2^{t+s}v-1$. Such a choice exists because the two values of
$2^{t+s}v-1$ differ by a nonzero residue modulo $3$. The congruence
ensures that $v3^{t+s}-1$ has \emph{exact} $2$-adic valuation $r$;
therefore $m$ is a positive odd integer for large $t$.

For $0\le j\le t$, induction gives
\[
 T^j(n)=2^{t+s-j}v3^j-1.
\]
Every intermediate starting value is odd, proving $T^t(n)=a$.
Likewise, for $0\le i\le s$,
\[
 T^i(a)=2^{s-i}v3^{t+i}-1.
\]
The first $s$ states are odd; each of the corresponding odd
steps strictly increases the state.  After these $s$ steps the
value is $v3^{t+s}-1=2^rm$, followed by $r$ consecutive halvings.

The bounds $v<2^{r+2}$ and $s=o(t)$ yield
\[
 \log X\le t\log3+(s+r+2)\log2,
 \qquad
 \log(m/X)\le s\log(3/2)-r\log2+O(1).
\]
Since $r/t\to1/45$ and
\[
 \frac{\log2}{45}>0.01\left(\log3+\frac{\log2}{45}\right),
\]
we obtain $m<X^{0.99}$ for all sufficiently large $t$.
Furthermore, $s\to\infty$, $s=o(t)$, $\log X=O(t)$, and
$\log\log X\le\log t+O(1)$; hence $b_X\le s\le t_X$ eventually.
The first $s$ steps from $a$ are odd and strictly increasing, so
$T^i(a)\ge a=X>X_1$ for all $0\le i<s$. Thus the prefix $s$
certifies membership in the live danger set. Because $s\ge2$,
$3a+1=2(2^{s-1}v3^{t+1}-1)$ with an odd factor in parentheses,
so $\nu_2(3a+1)=1$.

Finally, the $s$ initial odd steps move strictly above $a=X$.
The ensuing pure halving chain begins above $X$ and reaches
$m<X_1<X/2$. It therefore crosses $(X/2,X]$ \emph{exactly once},
with all its subsequent states at most $X/2$. This crossing occurs
before the first passage through $X_1$. Together with the initial
state $a$, there are exactly two shell visits.
\end{proof}

\section{Potential spikes: a strictly conditional observation}
For $y\in\N$ set $P(y)=\{w:\exists k\ge0,\ T^k(w)=y\}$ and let
$h(y)$ be its logarithmic density, assuming existence as in the
consolidated manuscript. Set
\[
 h_X^\downarrow(y)=h(y)\,\one_{\{\tau_{X_1}(y)<\infty\}},
 \qquad g_X(y)=y h_X^\downarrow(y).
\]
\begin{corollary}[Conditional polynomial live-set spikes]\label{cor:UF}
Assume, in addition to the established existence of $h$, the
\emph{external, not independently verified} hypothesis
\[
 \tag{UF}\label{eq:UF}
 \exists\kappa>0\quad\forall z\in\N\quad
 (3\nmid z\implies h(z)\ge\kappa/z).
\]
Define
\[
 \delta_*:=\frac{\log(3/2)}{\log3+(\log2)/45}=0.3639\ldots.
\]
For every fixed $0<\delta<\delta_*$ and arbitrarily large
$X$, there is an odd $a=X\in\calC^{\mathrm{live}}_{X,b_X}$
with $\nu_2(3a+1)=1$, $N_{(X/2,X]}^+(a)=1$, and
\[
 \boxed{\quad g_X(a)\ge X^\delta.\quad}
\]
Consequently, under (UF), a uniform odd-target $L^p$ bound in
$W_X$ fails for each $p>1/\delta_*=2.7475\ldots$, even with
the average restricted to the live parity danger set.
\end{corollary}
\begin{proof}
Since $T^t(n)=a$, $P(n)\subseteq P(a)$, so
$h(a)\ge h(n)\ge\kappa/n$. The stopping property of
Theorem~\ref{thm:family} gives $g_X(a)=ah(a)$. Also
\[
 \frac an\ge\frac12(3/2)^t,
 \qquad
 \log X\le t\bigl(\log3+(\log2)/45+o(1)\bigr).
\]
It follows that $g_X(a)\ge X^{\delta_*-o(1)}$.
For any $\delta<\delta_*$ this exceeds $X^\delta$ eventually.
The single term $y=a=X$ contributes at least $X^{p\delta-1}$ to
$X^{-1}\sum_{y\in W_X\cap\calC^{\mathrm{live}}_{X,b_X}}g_X(y)^p$.
Choose $\delta<\delta_*$ with $p\delta>1$.
\end{proof}
\begin{remark}[Evidence boundary]
Theorem~\ref{thm:family} is unconditional. Its conclusions do
\emph{not} require (UF), and its examples may have negligible
actual predecessor density. The polynomial-height claim in
Corollary~\ref{cor:UF} is conditional on an outside lower-bound
claim that has not been audited for this note. In particular,
no violation of the clipping-tail asymptotic is inferred: each
constructed orbit visits its own shell only twice.
\end{remark}

\section{The exact repair: certify actual later returns}
Fix a finite subprobability entrance law $\eta$ supported on states
whose orbits reach $[1,X_1]$. In a shell $I\subset (X_1,\infty)$,
let $\nu_n(y)$ be the $\eta$-mass whose $n$th visit to $I$ (before
first passage) is at $y\in I$, and let $s(n)=\sum_{y\in I}\nu_n(y)$.
Let
\[
 A_{I,b}:=\{y\in I:N_I^+(y)\ge b\}.
\]
\begin{proposition}[Exact return-certificate identity]\label{prop:certificate}
For every $n,b\ge1$,
\[
 \boxed{s(n+b)=\sum_{y\in A_{I,b}}\nu_n(y).}
\]
For every integer $K\ge0$,
\[
 \sum_{n>K}s(n)
 =\sum_{n=K+1}^{K+b}s(n)
 +\sum_{n>K}\nu_n(A_{I,b}).
\]
In particular, if $s(n)\le|\eta|$, then
\[
 \sum_{n>K}s(n)\le b|\eta|+\sum_{n\ge1}\nu_n(A_{I,b}).
\]
If, in addition, $\eta$ is the logarithmic first-excess entrance
law in the consolidated manuscript, all trajectories in its support
reach the cutoff, and $h^\downarrow_X(y)$ dominates the total
$\eta$-mass of starting values whose trajectory visits $y$, then
\[
 \sum_{n>K}s(n)\le b|\eta|+\sum_{y\in A_{I,b}}h_X^\downarrow(y).
\]
\end{proposition}
\begin{proof}
On the event of an $n$th visit at $y$, the future orbit is
\emph{deterministic}. It contains an $(n+b)$th visit precisely when
$N_I^+(y)\ge b$. This proves the first identity. Shift indices in
$\sum_{n>K+b}s(n)$ to obtain the second; bound the first $b$
terms by $b|\eta|$. A trajectory reaching the cutoff cannot revisit
the same $y>X_1$ before first passage, as that would trap it in a
cycle above $X_1$. Hence the events ``the $n$th visit is at $y$''
are disjoint in $n$, and
$\sum_{n\ge1}\nu_n(y)$ is at most the total entrance-source mass
whose orbit visits $y$. The stated predecessor-potential domination
proves the last inequality.
\end{proof}

\begin{corollary}[Return-certified shell-potential bound]
Under the entrance-law and first-passage hypotheses of
Proposition~\ref{prop:certificate}, suppose the weighted shell
excess of the consolidated manuscript is
\[
 \calE_X(L_X;K_X)=\sum_j2^{-j}\sum_{n>K_{X,j}}s_{X,j}(n),
 \quad I_{X,j}=(X2^{-j-1},X2^{-j}]\cap W_X,
\]
with $|\eta_{X,L_X}|=\mu_X(L_X)$. Then, for every $b\ge1$,
\[
 \boxed{\quad
 \calE_X(L_X;K_X)\le2b\,\mu_X(L_X)
 +\frac2X\sum_j\sum_{y\in A_{I_{X,j},b}}
 y\,h_X^\downarrow(y).\quad}
\]
For $b\ge2$, the targets of Theorem~\ref{thm:family} are excluded
from the return-certified set for their own shell, although they
belong to the corresponding live danger set.
\end{corollary}
\begin{proof}
Apply Proposition~\ref{prop:certificate} to each shell. Sum
with weights $2^{-j}$. Since $\sum_j2^{-j}\le2$ and
$2^{-j}<2y/X$ for $y\in I_{X,j}$, the displayed bound follows.
The final assertion is Theorem~\ref{thm:family}(3).
\end{proof}

\section{Interpretation: what is \emph{not} yet established}
The return-certified estimate is strictly more selective than the
survival-gated parity-cylinder estimate. Its condition is, however,
\emph{dynamical}: the set $A_{I,b}$ requires actual later returns,
not just a combinatorial parity word. No decay estimate for
$X^{-1}\sum_{j,y\in A_{I_{X,j},b}}y h_X^\downarrow(y)$ is established
here, and such an estimate must not be inferred from the cardinality
of parity cylinders. The key open problem remains a \emph{weighted}
upper bound for genuine multi-return states under the logarithmic
first-excess measure.  This is not a proof of $T_X=o(X)$ or a new
Lean kernel certificate. Nothing in the consolidated paper has been
modified.

\begin{thebibliography}{9}
\bibitem{Tao}
T. Tao, \emph{Almost all orbits of the Collatz map attain almost
bounded values}, Forum Math. Pi \textbf{10} (2022), e12.
\bibitem{PriorNote}
Internal note, \emph{A Last-Mile Audit: Polynomial Predecessor
Spikes and Survival-Gated Dangerous Cylinders}, 17 November 2026.
\bibitem{UFclaim}
Fourier\_wren, \emph{A uniform linear lower bound for Collatz
predecessor counts, from joint Syracuse atoms at fixed valuation},
2026 online manuscript, \url{https://rstation.ai/problems/collatz/papers/10/v2}.
The asserted uniform lower bound is unverified here.
\end{thebibliography}
\end{document}
